\chapter{Конечная конфигурация с нулём Хиггса: солитон или сфалерон}
\label{ch:v6-spectral-transition-higgs-zero-finite-energy-saddle-gate}

\section{Вопрос после топологического запрета}

Предыдущая глава установила, что вакуум одного стандартного хиггсовского
дублета равен $S^3$ и не имеет классов $\pi_0$, $\pi_1$ или $\pi_2$,
которые принуждали бы локальное ядро $H=0$. Это не исключает
нетопологической стационарной конфигурации. Настоящий гейт проверяет два
разных уровня:
\begin{enumerate}
 \item чистый канонический скалярный сектор Хиггса;
 \item полный электрослабый gauge--Higgs сектор.
\end{enumerate}

Цель состоит не в воспроизведении наблюдаемой энергии сфалерона, а в
классификации физического типа возможного нуля $H$ и связанной смены ранга
\begin{equation}
 W_\nu(H)=\widetilde H\widetilde H^\dagger.
\end{equation}

\section{Чистый скалярный сектор и тождество Деррика}

После вычитания вакуумной энергии статический функционал одного дублета
имеет вид
\begin{equation}
 E_H[H]=T_H+V_H
 =\int_{\mathbb R^3}\left(
 |\nabla H|^2+
 \lambda\left(H^\dagger H-\frac{v^2}{2}\right)^2
 \right)d^3x,
 \label{eq:v6-higgs-scalar-energy}
\end{equation}
где $T_H,V_H\geq0$. Для масштабированного поля
$H_R(x)=H(x/R)$ получаем
\begin{equation}
 E_H(R)=R T_H+R^3V_H.
 \label{eq:v6-higgs-derrick-scaling}
\end{equation}
Поэтому
\begin{equation}
 \left.\frac{dE_H}{dR}\right|_{R=1}=T_H+3V_H.
\end{equation}
У нетривиальной локализованной конфигурации правая часть положительна и
не может удовлетворять условию стационарности. Это стандартный
масштабный запрет Деррика \cite{higgs-saddle-derrick1964}.

В частности, принудительно заданная депрессия $|H(0)|=0$ может иметь
конечную пробную энергию, но при свободной вариации либо сжимается, либо
размывается к однородному вакууму. Она не является статическим солитоном.

\section{Регулярность фиксированного радиального направления}

Для анзаца $H(x)=f(r)H_0$ с постоянным внутренним направлением уравнение
имеет структуру
\begin{equation}
 f''+\frac2r f'=\lambda f(f^2-v^2).
 \label{eq:v6-higgs-radial-equation}
\end{equation}
Гладкость скалярной функции $f(|x|)$ в начале требует $f'(0)=0$.
Если дополнительно потребовать $f(0)=0$, разложение в ряд сначала удаляет
линейный коэффициент из-за члена $2f'/r$, затем последовательно удаляет
все остальные коэффициенты. Регулярная локальная ветвь есть $f\equiv0$ и
не может прийти к $f(\infty)=v$.

Разрешение поведения $f\sim r$ требует компенсирующего углового
внутреннего направления, как у ежа. Но один дублет не имеет
нетривиального $\pi_2(S^3)$, а фиксированный $f(r)H_0$ с линейным началом
не является гладким полем в декартовых координатах.

\begin{theorem}[Нет чисто хиггсовского статического комка]
Канонический двухпроизводный функционал одного комплексного дублета в
трёх пространственных измерениях не допускает нетривиального устойчивого
статического солитона конечной энергии. В частности, локальный нуль $H$
не стабилизирует спектральную опору $W_\nu$ без дополнительного поля,
заряда, времени или высшего производного члена.
\end{theorem}

\section{Почему gauge-поле меняет масштабный баланс}

Для полного электрослабого статического функционала обозначим через
$E_F$ энергию калибровочной кривизны, через $E_D$ ковариантную кинетику
Хиггса и через $E_V$ потенциал. При согласованном масштабировании
\begin{equation}
 A_i^{(R)}(x)=R^{-1}A_i(x/R),
 \qquad H_R(x)=H(x/R)
\end{equation}
получается
\begin{equation}
 E_{\rm EW}(R)=R^{-1}E_F+R E_D+R^3E_V.
 \label{eq:v6-electroweak-scaling}
\end{equation}
Вириальное условие
\begin{equation}
 -E_F+E_D+3E_V=0
 \label{eq:v6-electroweak-virial}
\end{equation}
уже не противоречит положительности: gauge-кривизна способна балансировать
скалярные члены. Поэтому масштабный запрет чистого Хиггса не запрещает
конечную стационарную gauge--Higgs конфигурацию.

\section{Электрослабый сфалерон}

Именно такой конфигурацией является сфалерон Клинкхамера--Мэнтона ---
конечное статическое решение уравнений Вайнберга--Салама, расположенное на
вершине барьера между вакуумами с различными числами Черна--Саймонса
\cite{higgs-saddle-klinkhamer-manton1984}. Его стандартный хиггсовский
профиль удовлетворяет
\begin{equation}
 H(0)=0,
 \qquad |H(r)|\longrightarrow\frac{v}{\sqrt2}
 \quad(r\to\infty).
\end{equation}
Следовательно, он действительно реализует пространственный переход
\begin{equation}
 \rank W_\nu(0)=0,
 \qquad
 \rank W_\nu(r)=1\quad(r>0)
\end{equation}
для стандартного профиля.

Но сфалерон является не минимумом, а седлом. Современный полный расчёт со
стандартными параметрами снова находит единственную отрицательную моду
\cite{higgs-saddle-matchev-verner2025}. Поэтому малое возмущение ведёт к
одному из соседних вакуумов; локальный нуль ранга не сохраняется как
частица.

\begin{proposition}[Смена ранга как переходное событие]
Полный электрослабый gauge--Higgs сектор допускает конечную стационарную
конфигурацию с локальным $H=0$, но она имеет отрицательную моду. Поэтому
сфалерон реализует смену ранга спектральной опоры как событие перехода
между вакуумами, а не как устойчивый материальный объект.
\end{proposition}

\section{Что наследует и чего не наследует проект}

Проект уже содержит стандартные представления
$SU(2)_L\times U(1)_Y$, хиггсовский дублет и положительную
самосвязь $\lambda_H$. Этого достаточно для структурного чтения
сфалеронного маршрута. Однако абсолютные $g_2$, $g_Y$, $v$ и полный
низкоэнергетический перенос нормировок в проекте имеют разные статусы.
Поэтому численное значение энергии сфалерона не объявляется новым
предсказанием S2T.

Также нельзя отождествлять одну отрицательную моду сфалерона с виртуальной
частицей. В строгом языке это направление распада классического седла.
Квантовое чтение требует оператора фермионного спектрального потока и
аномального изменения заряда.

\begin{nogocriterion}[Седло не является материей]
Наличие конечной энергии, локального $H=0$ и фермионной нулевой моды в
одной критической конфигурации недостаточно для идентификации частицы.
Необходим либо положительный физический гессиан после удаления симметрий,
либо иной сохраняемый квантовый сектор, запрещающий распад.
\end{nogocriterion}

\section{Следующий различающий тест}

Отрицательная мода закрывает сфалерон как устойчивый комок, но открывает
естественный динамический вопрос. Следующий гейт
\path{version6_spectral_transition_sphaleron_spectral_flow_gate}
должен сопоставить:
\begin{enumerate}
 \item изменение числа Черна--Саймонса на единицу;
 \item спектральный поток фермионного оператора через ноль;
 \item пакет одного поколения $H_{15}$ и его разложение $12+3$;
 \item Real/KO6-удвоение и тёплицев класс пятнадцать.
\end{enumerate}

Нулевая гипотеза: стандартный аномальный спектральный поток видит
хиральные электрослабые дублеты и не совпадает с тёплицевым классом
$15$ как с числом локализованных частиц. Совпадение одного целого без
явной операторной карты не считается мостом.

Стоп-критерий: если коэффициент аномалии, число пересечений и Real-класс
не могут быть согласованы на одном операторе без добавленных состояний,
сфалеронный маршрут не объясняет материальный пакет проекта и остаётся
только стандартным электрослабым переходом.

\begin{protocol}
\begin{enumerate}
 \item чистый скалярный масштабный баланс: \textbf{невозможен};
 \item гладкий фиксированно-направленный профиль с $H(0)=0$:
 \textbf{тривиален};
 \item gauge--Higgs вириальный баланс: \textbf{допустим};
 \item конечная конфигурация с $H(0)=0$: \textbf{сфалерон существует};
 \item ранг $W_\nu$ в центре сфалерона: \textbf{ноль};
 \item физический гессиан сфалерона: \textbf{имеет отрицательную моду};
 \item стабильная частица из одного Хиггса: \textbf{не получена};
 \item смена ранга как переходное событие: \textbf{поддержана};
 \item следующий тест: \textbf{фермионный спектральный поток};
 \item физическое замыкание: \textbf{не достигнуто}.
\end{enumerate}
\end{protocol}

Машинный сертификат сохранён в
\path{s2t_v6_spectral_transition_higgs_zero_finite_energy_saddle_gate_results.json}.

\begin{thebibliography}{9}
\bibitem{higgs-saddle-derrick1964}
G.~H.~Derrick,
\emph{Comments on Nonlinear Wave Equations as Models for Elementary
Particles},
Journal of Mathematical Physics 5 (1964), 1252--1254.

\bibitem{higgs-saddle-klinkhamer-manton1984}
F.~R.~Klinkhamer, N.~S.~Manton,
\emph{A Saddle-Point Solution in the Weinberg--Salam Theory},
Physical Review D 30 (1984), 2212--2220.

\bibitem{higgs-saddle-matchev-verner2025}
K.~T.~Matchev, S.~Verner,
\emph{The Electroweak Sphaleron Revisited: I. Static Solutions, Energy
Barrier, and Unstable Modes},
arXiv:2505.05607.

\bibitem{higgs-saddle-klinkhamer2003}
F.~R.~Klinkhamer,
\emph{Sphalerons, Spectral Flow, and Anomalies},
arXiv:hep-th/0304167.
\end{thebibliography}